%% Generated by tools/build_guide.py from UNIT-REGISTRY.md.
%% Do not edit by hand: edit the source Markdown or the registry.
\chapter{语言基础}
\guidechaptermark{语言基础}
\label{ch:language-basics}

\guidechapterintro{本章通过简短的示例，帮助你读懂并编写 Croqtile 程序。先从完整程序的结构和各类函数的分工看起，再逐步加入形状、迭代、分块选择、数据搬运与协调。每个示例都说明新构造在程序中的作用，让你在具体的计算过程中理解它的写法。}

\section{Croqtile 程序的两个部分}
\label{sec:2-1}

一个 \co@.co@ 文件同时包含普通 C++ 代码和 Croqtile 代码。主机程序准备数据、调用核函数并使用结果；带 \co@__co__@ 前缀的函数描述计算与数据搬运，由 Croqtile 编译器翻译成目标代码。


\subsection*{主机、Tileflow 与设备代码}

普通 C++ 代码还可以包含设备辅助函数。因此，完整示例中常见三个区域：设备辅助函数、调用它的 \co@__co__@ 函数，以及主机侧的 \co@main()@。下面按这个顺序列出它们。

\begin{lstlisting}[language=croqtile,caption={单份源文件中的三个区域}]
__device__ void device_function(...) {
  // High-performance device kernel implementation
}
// Tileflow program: orchestrating data movement
__co__ void croqtile_function(...) {
  // ... croqtile code ...
  device_function(...);
  // ...
}
// Host program: typically runs on CPU
void main() {
  // ... prepare data ...
  croqtile_function(...);
  // ...
}
\end{lstlisting}

程序从主机侧开始执行。主机代码使用普通 C++，负责准备输入、调用 Croqtile 函数并处理结果。

\co@__co__@ 函数构成 Tileflow 程序，负责组织并行工作和数据搬运。它将数据放到合适的缓冲区，安排设备端计算，并把结果交回调用方。

\co@__device__@ 辅助函数执行设备端计算。Croqtile 编译器主要翻译 \co@__co__@ 区域，普通 C++ 的主机代码和设备代码交给后续工具链处理。\textless{}br\textgreater{}\textless{}br\textgreater{}\textbf{提示：} \co@__device__@ 主要用于在必要时复用特定硬件的函数，简化开发。对于 Croqtile 已提供较完善支持的硬件，通常直接使用语言提供的计算与数据搬运能力即可，无需通过设备辅助函数自行补充底层实现。



\section{逐行阅读核函数}
\label{sec:2-2}

\subsection*{Croqtile 函数}

先看函数签名。普通 C++ 函数的声明主要约定类型，Croqtile 还把数据形状写在参数和返回值中。这样，读者可以在进入函数体之前，先了解输入内容与预期的结果形状。

\begin{lstlisting}[language=croqtile,caption={核函数签名}]
__co__ s32 [4, 8] ele_add(s32 [4, 8] lhs, s32 [4, 8] rhs) {
\end{lstlisting}

前缀 \co@__co__@ 标记 Croqtile 函数。\co@s32 [4, 8]@ 把 32 位有符号整数元素与二维形状组合起来，表示四行八列。输入与返回值都采用 \co@type [shape] name@ 的形式。编译器据此检查秩并推导索引范围；依赖调用方实际尺寸的约束，还可能需要运行时检查。

常见元素类型包括 \co@s32@、\co@f16@、\co@bf16@、\co@f32@，以及 \co@f8_e4m3@ 和 \co@f8_e5m2@。本例使用 \co@s32@，便于观察逐元素加法的结果。

\begin{lstlisting}[language=croqtile,caption={分配输出}]
s32 [lhs.span] output;
\end{lstlisting}

这一行声明输出张量。\co@lhs.span@ 取出输入的完整形状，因此 \co@output@ 也具有 \co@[4, 8]@ 的形状。修改输入形状时，这条声明的输出形状会随之改变；本例中固定的函数签名和 \co@parallel@ 范围仍需同步调整。

\begin{lstlisting}[language=croqtile,caption={并行函数体}]
parallel {i, j} by [4, 8]
  output.at(i, j) = lhs.at(i, j) + rhs.at(i, j);
\end{lstlisting}

语句 \co@parallel {i, j} by [4, 8]@ 声明 32 个独立实例，每个实例负责一个元素位置。例如 \co@(i, j) = (1, 3)@ 的实例读取两份输入的这一位置，再把和写入 \co@output@ 的相同位置。不同实例的写入位置不重叠，也不依赖其他实例产生的结果，因此这些工作可以并行执行。

没有空间说明符时，由编译器选择实例到执行资源的映射。语句确定工作彼此独立，并不要求每个元素都对应一个物理线程，也不要求所有实例同时运行。

\begin{lstlisting}[language=croqtile,caption={返回结果}]
return output;
\end{lstlisting}

返回结果的形状与元素类型必须符合函数签名。主机侧收到的 \co@croq::spanned_data@ 拥有结果缓冲区，并携带形状信息，可通过 \co@[i][j]@ 访问元素。


\subsection*{\texttt{\_\_co\_\_ auto} 返回类型}

若结果形状更适合在函数体中描述，可以使用 \co@auto@。对于 \co@__co__ auto matmul(...)@，编译器从 \co@return output@ 所返回的缓冲区推导类型。返回值仍然有元素类型和形状，\co@auto@ 只是省去了签名中重复的声明。若希望直接在接口中约定输出形状，则可以写出完整返回类型。

逐元素算术、标量矩阵乘法以及返回缓冲区的 MMA 核函数，都遵循相同的返回约定。函数内部怎样计算元素，并不影响这一规则。



\section{主机侧 API：输入与输出}
\label{sec:2-3}

\subsection*{主机程序}

主机侧提供数据存储和形状信息。加法示例先创建两份拥有存储的缓冲区，填入数据，再将视图传给核函数。主机保留输入的所有权，调用结束后还可以用原始数据检查结果。

\begin{lstlisting}[language=croqtile,caption={创建带形状的输入}]
auto lhs = croq::make_spandata<croq::s32>(4, 8);
auto rhs = croq::make_spandata<croq::s32>(4, 8);
lhs.fill_random(-10, 10);
rhs.fill_random(-10, 10);
\end{lstlisting}

调用 \co@croq::make_spandata<T>(dims...)@ 分配由对象拥有的主机缓冲区。模板参数选择元素类型，实参给出各维大小。随后调用 \co@fill_random@，在指定范围内初始化元素。

\begin{lstlisting}[language=croqtile,caption={调用核函数}]
auto res = ele_add(lhs.view(), rhs.view());
\end{lstlisting}

主机侧的调用形式与普通 C++ 函数调用相同。\co@.view()@ 从 \co@spanned_data@ 创建不拥有缓冲区的 \co@spanned_view@，供核函数访问输入；返回的 \co@res@ 则拥有结果缓冲区。

\co@main@ 最后逐个检查结果元素，验证核函数的计算是否正确。这个检查沿用相同的 \co@(i, j)@ 位置，把两个原始输入的和与返回结果比较，从主机侧完成一次完整的准备输入、调用与验证的过程。


\subsection*{主机代码}

下面的主机代码使用相同的 \co@make_spandata@ 和 \co@.view()@ 调用方式。

\begin{lstlisting}[language=croqtile,caption={主机代码}]
int main() {
  auto lhs = croq::make_spandata<croq::s32>(128, 256);
  auto rhs = croq::make_spandata<croq::s32>(256, 256);
  lhs.fill_random(-10, 10);
  rhs.fill_random(-10, 10);
  auto res = matmul(lhs.view(), rhs.view());
  for (size_t i = 0; i < res.shape()[0]; ++i)
    for (size_t j = 0; j < res.shape()[1]; ++j) {
      int ref = 0;
      for (size_t k = 0; k < lhs.shape()[1]; ++k)
        ref += lhs[i][k] * rhs[k][j];
      croq_assert(ref == res[i][j], "values are not equal.");
    }
  std::cout << "Test Passed\n" << std::endl;
}
\end{lstlisting}

调用后，程序将每个结果元素与 C++ 参考计算比较；结果不一致时，断言会报错。


\subsection*{Croqtile 与 C++ 函数之间的参数传递}

C++ 指针需要与维度信息一起包装成带形状的数据，再传入 Croqtile 函数。相关工具由 \co@choreo.h@ 提供。

\begin{lstlisting}[language=croqtile,caption={Croqtile 与 C 函数之间的参数传递}]
#include "choreo.h"
void bar(const float* data, unsigned size) {}
__co__ void foo(f32 mdspan<2> d) {
  parallel p by 6 {
    call bar(d, |d|);     // Call the C++ function
  }
}
void foobar(float* a) {
  foo(croq::make_spanview<2>(a, {1, 2}));
}
\end{lstlisting}

\co@make_spanview@ 创建输入视图，不复制底层数据。示例中，\co@foo@ 的参数 \co@d@ 使用 \co@f32 mdspan<2>@ 描述形状；调用 \co@bar@ 时，可以通过隐式转换传递底层数据，\co@|d|@ 则给出元素总数。

标量值也可以在 C++ 与 Croqtile 之间传递；\co@ituple@ 仅用于 Croqtile 函数内部。








\subsection*{输入参数}

Croqtile 函数的输入参数包括带形状数据，以及标量：整数、布尔值或浮点元素类型。

参数可以具名，也可以省略名字：

\begin{lstlisting}[language=croqtile,caption={输入参数}]
__co__ void foo(f32 [7] a, int) {...}
\end{lstlisting}

本例的第二个参数没有名字，因此函数体无法引用它，但声明仍然合法。


\subsection*{C++ 主机：带形状的数据实参}

普通指针或数组只描述存储，Croqtile 的带形状数据还记录各维度大小。调用前，需要把数据地址与形状组合起来。

可以使用 \co@choreo.h@ 中的包装工具。Croqtile 编译环境会自动包含这个头文件。

\begin{lstlisting}[language=croqtile,caption={传入带形状的缓冲区}]
__co__ void foo(f32 [7, 16] input, f32 [7, 16] output) {...}
void entry(float * input, float* output) {
  // foo(input, output); <-- result in error since missing shape
  foo(croq::make_spanview<2>(input, {7, 16}),
      croq::make_spanview<2>(output, {7, 16}));
}
\end{lstlisting}

\co@croq::make_spanview@ 以秩(rank)作为模板实参，接收数据指针和维度列表，生成 \co@croq::spanned_view@。它保存指针与形状，不拥有底层缓冲区。简化定义如下。

\begin{lstlisting}[language=croqtile,caption={主机侧缓冲区包装器}]
template <typename T, size_t Rank>
class spanned_view {
  T* ptr;                 // typed pointer to the data
  size_t dims[Rank];      // ranked dimensions
};
\end{lstlisting}

维度也可以由运行时值提供：

\begin{lstlisting}[language=croqtile,caption={由运行时维度构造的调用}]
__co__ void foo(f32 [7, 16] input) {...}
void entry(float * input, float* output, int M, int N) {
  foo(croq::make_spanview<2>(input, {M, N}));
}
\end{lstlisting}

本例的 \co@input@ 使用运行时形状。进入 Tileflow 程序时，系统检查它是否满足 \co@M == 7@ 和 \co@N == 16@；不满足时会终止执行，以免继续使用不成立的形状假设。

输入的带形状数据引用原有缓冲区。把这个包装对象传入函数，不会复制整块输入数据。





\subsection*{返回类型与类型推导}

返回值和输入一样，可以是带形状数据、标量，或者不返回结果（\co@void@）。

除了显式声明返回类型，也可以让编译器推断：

\begin{lstlisting}[language=croqtile,caption={返回类型与类型推导}]
__co__ auto foo() {
  f32 [7, 16] result;
  // ...
  return result;
}
\end{lstlisting}

这里，编译器根据 \co@result@ 将 \co@auto@ 推导为 \co@f32 [7, 16]@。当输出形状在函数体内推导时，这种写法可以避免修改推导式后还要同步修改返回声明。若固定返回形状是接口的一项要求，则适合显式声明类型。


\subsection*{C++ 主机：形状返回}

主机侧收到的结果拥有自己的缓冲区，具有值语义，便于在调用结束后继续使用。

\begin{lstlisting}[language=croqtile,caption={返回带形状的缓冲区}]
__co__ auto foo() {
  f32 [7, 16] output;
  // ...
  return output;  // semantically copy
}
void entry() {
  auto result = foo(); // move the 'spanned_data' to caller
  // ...
}
\end{lstlisting}

示例中，\co@output@ 作为结果返回，主机侧类型为 \co@croq::spanned_data@。实现通过返回值优化避免复制整块缓冲区。简化定义如下。

\begin{lstlisting}[language=croqtile,caption={主机侧结果包装器}]
template <typename T, size_t Rank>
class spanned_data {
  std::unique_ptr<T[]> ptr; // move only pointers
  size_t dims[Rank];        // ranked dimensions
};
\end{lstlisting}

\co@spanned_view@ 只引用数据，\co@spanned_data@ 则管理缓冲区的所有权。返回值优化允许在交付结果时省去数据复制。

主机侧可以用 \co@.rank()@、\co@.shape()@ 查询结果形状，而不必依赖核函数内部的声明。



\section{形状与 \texttt{mdspan}}
\label{sec:2-4}

\subsection*{\texttt{span} 与 \texttt{span(i)}}

成员 \co@tensor.span@ 返回完整形状，\co@tensor.span(i)@ 返回第 i 维的大小。矩阵乘法的输出行数来自左输入，列数来自右输入：

\begin{lstlisting}[language=croqtile,caption={span 与 span(i)}]
s32 [lhs.span(0), rhs.span(1)] output;
\end{lstlisting}

这个例子从两个不同操作数取得结果尺寸，使输出声明与输入保持关联。需要匹配已搬运分块形状的暂存缓冲区，也可以用同样的方法声明。

\begin{figure}[htbp]\centering
\includegraphics[width=\linewidth,height=0.56\textheight,keepaspectratio]{f04-fig_span_light.png}
\caption{对 [4, 8] 张量，.span 返回完整形状 [4, 8]；.span(0) 和 .span(1) 分别返回 4 和 8。}
\end{figure}


形状本身也可以命名并参与计算。下面先声明一个形状，再说明怎样从它推导出新的形状。


\subsection*{用 \texttt{mdspan} 定义形状}

关键字 \co@mdspan@ 表示多维范围，用来描述形状各维的大小，本身不分配元素存储。例如 \co@[7, 8]@ 的秩为 2，描述 56 个位置；这些位置采用 \co@f32@ 还是 \co@s32@ 元素，则是另一项选择。

显式声明的写法如下：

\begin{lstlisting}[language=croqtile,caption={声明形状}]
mdspan s0 : [7, 8]; // Defines a 2D shape with dimensions [7, 8]
mdspan<1> s1 : [3]; // Defines a 1D shape with dimensions [3]
\end{lstlisting}

\co@:@ 引出初始化表达式，\co@[]@ 中依次列出各维度大小。\co@s0@ 的形状是 \co@[7, 8]@，\co@s1@ 的形状是 \co@[3]@。

\co@mdspan@ 后的尖括号可以指定秩，也就是维度数量。它必须与初始化表达式一致，否则编译器会报错。

\begin{lstlisting}[language=croqtile,caption={秩不匹配}]
mdspan<3> s2 : [64, 32]; // error: the rank of mdspan is inconsistent
\end{lstlisting}

编译器也可以从初始化表达式推断类型，此时省略 \co@mdspan@：

\begin{lstlisting}[language=croqtile,caption={省略 mdspan 关键字}]
s3 : [7, 8, 9]
\end{lstlisting}

\co@s3@ 的秩为 3，形状为 \co@[7, 8, 9]@。\co@mdspan@ 必须在声明时初始化；省略类型只简化写法，不改变这项要求。


\subsection*{导出 \texttt{mdspan}}

分块、填充和重塑都需要从已有形状推导新形状。例如，把每维大小除以分块因子，可以得到每块形状；给某维增加长度，则能描述填充后的区域。把这些关系写成形状表达式，修改原始尺寸时就不必逐处更新数字。

可以对维度做算术运算，再组合成新的形状：

\begin{lstlisting}[language=croqtile,caption={原始形状}]
shape : [128, 64]; // initial shape
new-shape0 : shape [(0) / 2, (1) / 4, 1];  // tile and reshape: [1, 64, 16]
new-shape1 : shape [(1) + 2, (0) / 16];    // pad and reshape: [66, 8]
\end{lstlisting}

\co@new-shape0@ 将原形状的两个维度分别除以 \co@2@ 和 \co@4@，并增加一个长度为 \co@1@ 的维度。

同一计算也可以展开写成：

\begin{lstlisting}[language=croqtile,caption={导出的形状}]
new-shape0: [shape(0) / 2, shape(1) / 4, 1];
\end{lstlisting}

展开后，\co@shape(0)@ 和 \co@shape(1)@ 分别取出原形状的两个维度。两种写法表达的是同一组计算。

\co@new-shape1@ 将原形状的第二个维度加 \co@2@，把它放到新形状的第一维；新形状的第二维由原形状的第一维除以 \co@16@ 得到。

也可以直接对整个形状做运算或拼接：

\begin{lstlisting}[language=croqtile,caption={重新绑定形状}]
shape : [32, 72]
new-shape0 : shape;  // [32, 72]
new-shape1 : shape + 1;  // [33, 73]
new-shape2 : shape / 4;  // [8, 18]
new-shape3 : [shape, 6]; // [32, 72, 6]
\end{lstlisting}

对整个形状进行算术运算时，运算逐维应用。例如 \co@[128, 64]@ 加 1 得到 \co@[129, 65]@，并不会增加一个维度。下面两种表达式写的是同一项计算：

\co@new-shape1 : shape + 1;@

可展开为：

\co@new-shape1 : shape [(0) + 1, (1) + 1];@

在初始化表达式中放入一个 \co@mdspan@，则会拼接它的各个维度。例如：

\co@new-shape3 : [shape, 6];@

可展开为：

\co@new-shape3 : shape [(0), (1), 6];@

这些简写能保留形状之间的关系，减少重复列出维度的代码。


\subsection*{带匿名维度的 \texttt{mdspan}}

前面的形状使用固定数字。如果某个维度要由调用方在运行时提供，可以把对应位置写成 \co@?@，保留其他已知维度：

\begin{lstlisting}[language=croqtile,caption={含未知维度的核函数}]
__co__ auto foo(s32 [?, 1, 2] input) { ... }
\end{lstlisting}

\co@?@ 表示一个运行时提供的正整数维度：

\begin{lstlisting}[language=croqtile,caption={由调用方提供该维度}]
void bar(int * data, int dim0) {
  foo(croq::make_spanview<3>(data, {dim0, 1, 2}));
}
\end{lstlisting}

本例第一维的大小来自调用方，可以由配置或上层框架确定。

如果只需使用某份输入的一维大小，匿名维度已经足够。但当两份输入的某个维度必须相等时，它没有为这项关系提供共同的名字。矩阵乘法恰好需要这样的约束：左操作数的列数必须等于右操作数的行数。


\subsection*{更进一步：符号维度}

当多个参数需要引用同一个未知维度时，只写 \co@?@ 还不能直观表达它们的关系。符号维度给这些大小命名，让共享维度和形状约束出现在函数签名中：

\begin{lstlisting}[language=croqtile,caption={更进一步：符号维度}]
__co__ auto Matmul(s32 [M, K] lhs, s32 [K, N] rhs) {
  int tile_m = 32, tile_n = 8, tile_k = 16;
  s32 [M / tile_m, K / tile_k] tiled_lhs;
  s32 [K / tile_k, N / tile_n] tiled_rhs;
  // ...
}
\end{lstlisting}

本例用 \co@M@、\co@K@、\co@N@ 表达两个输入的形状。相同的名字明确了共享维度，便于阅读、维护和检查参数之间的关系。



\section{整数、整数元组(ituple)与有界值}
\label{sec:2-5}

\subsection*{标量类型}

用于程序控制的标量包括整数和布尔值。\co@int@ 是 32 位有符号整数，范围为 -2³¹ 到 2³¹-1，可用于声明变量和函数参数。

\begin{lstlisting}[language=croqtile,caption={标量类型}]
int a;
__co__ int foo(int b) { return b; }
\end{lstlisting}

整数支持 C++ 风格的算术、移位和逻辑运算。程序控制使用 \co@int@，不提供其他位宽或无符号的控制整数类型。

\co@bool@ 的运算以及它与整数之间的转换，采用 C++ 的对应规则。


\subsection*{整数元组类型}

整数元组按顺序保存多个整数，常用于多维索引。把行、列等坐标放在同一个元组中，可以整体传递或参与逐维计算；需要某个分量时，再按位置取出。元素顺序因此也是索引含义的一部分。

用花括号列出各个元素即可定义 \co@ituple@：

\begin{lstlisting}[language=croqtile,caption={整数元组类型}]
ituple index = {5, 4, 3, 2, 1};  // It defines a tuple of 5 elements
index = {a, b};  // 'a' and 'b' are existing integers
\end{lstlisting}


\subsection*{\texttt{mdspan} 与 \texttt{i-tuple} 上的运算}

\co@mdspan@ 可以与整数元组按维度运算。下面用一组分块因子计算新的形状。

\begin{lstlisting}[language=croqtile,caption={mdspan 与 i-tuple 上的运算}]
sp : [6, 8];
tiling_factor = {3, 2};
spn : sp / tiling_factor;   // spn is defined as [2, 4];
\end{lstlisting}

支持的运算包括：

\begin{itemize}[leftmargin=1.6em]
\item \co@mdspan@ / \co@i-tuple@
\item \co@mdspan@ + \co@i-tuple@
\item \co@mdspan@ \% \co@i-tuple@
\item \co@mdspan@ * \co@i-tuple@
\item \co@mdspan@ - \co@i-tuple@
\end{itemize}

元组写法将逐维计算集中在一个表达式中，更容易看出形状之间的关系。


整数描述单个索引或循环参数，整数元组描述多维索引。它们在声明时初始化，并可与形状一起用于后续检查。


\subsection*{用于程序控制的整数}

Tileflow 程序使用整数来控制循环和索引。为此提供的控制类型是 32 位有符号整数 \co@int@。

可以像 C++ 一样声明并初始化整数：

\begin{lstlisting}[language=croqtile,caption={已初始化的整数}]
int a = 1;
\end{lstlisting}

也可以省略类型，由初始化表达式推断：

\begin{lstlisting}[language=croqtile,caption={对整数赋值}]
a = 1;
\end{lstlisting}

一行中可以声明多个整数：

\begin{lstlisting}[language=croqtile,caption={一次定义多个整数}]
int a = 1, b = 2, c = 3;
d = 4, e = 5;
\end{lstlisting}

\co@mdspan@ 不能直接作为函数参数；整数则可以在主机、Tileflow 程序和设备函数之间传递。

\begin{lstlisting}[language=croqtile,caption={整数参数}]
__co__ void foo(int a) {
  parallel p by 1 {
    call kernel(a, 3);
  }
}
\end{lstlisting}

整数必须在声明时初始化，之后不能重新赋值：

\begin{lstlisting}[language=croqtile,caption={未初始化的声明}]
int a;  // error: declaration without initialization
int b = 1;
b = 3;  // error: re-assignment
\end{lstlisting}

这里的整数是不可变值。因此，声明中的 \co@=@ 用于初始化，而不是修改已有变量。





\subsection*{声明整数元组}

\co@ituple@ 将多个整数组合成一个元组，也称 \co@i-tuple@。显式声明如下：

\begin{lstlisting}[language=croqtile,caption={i-tuple 字面量}]
ituple a = {1, 2, 3};
b = {4, 5, 6};  // utilize the type inference
\end{lstlisting}

类型可以由初始化表达式推断。区分省略类型的 \co@ituple@ 和 \co@mdspan@ 时，注意以下两点：

\begin{itemize}[leftmargin=1.6em]
\item \co@ituple@ 的初始化表达式跟在赋值号 \co@=@ 之后（和整数一样），\co@mdspan@ 则在 \co@:@ 之后初始化；
\item \co@ituple@ 的初始化表达式用花括号括起，而不是 \co@mdspan@ 使用的 \co@[]@。
\end{itemize}

也可以显式指定秩，让编译器检查元素个数：

\begin{lstlisting}[language=croqtile,caption={秩不一致}]
ituple<3> a = {1, 2};  // error: inconsistent rank
\end{lstlisting}


\subsection*{元组运算与形状计算}

圆括号 \co@()@ 取出元组中的一个元素，下标从 \co@0@ 开始。也可以让整个元组参与计算，把几个坐标分量的变化写在同一表达式里：

\begin{lstlisting}[language=croqtile,caption={赋一个新元组}]
a = {3, 4};
b = {a(0), 1, a(1)};  // '()' to retrieve the element value
c = a {(0), (1), 2};  // syntax sugar
d = {a, 5, 6};        // concatenate
e = a + 1;            // addition is applied elementwise
\end{lstlisting}

\co@ituple@ 与整数、\co@mdspan@ 一样，必须初始化且不能重新赋值。

元组常用于逐维修改形状：

\begin{lstlisting}[language=croqtile,caption={用 i-tuple 索引 mdspan}]
shape : [7, 18, 28];
tiling_factors = {1, 2, 4};
tiled_shape : shape / tiling_factors;
padded_shape : shape + {2, 0, 2};
\end{lstlisting}

\co@shape / tiling_factors@ 按各维的分块因子计算 \co@tiled_shape@，\co@shape + {2, 0, 2}@ 则得到 \co@padded_shape@。按维度相加时，两者的秩必须一致。

\begin{lstlisting}[language=croqtile,caption={非法的形状表达式}]
shape : [7, 8, 9] + {1, 2}; // error: inconsistent rank
\end{lstlisting}


\subsection*{整数与元组的使用规则}

一个控制值用整数表示，一组有顺序的值用 \co@ituple@ 表示。形状 \co@[7, 14]@ 描述范围，元组 \co@{1, 2}@ 则可以表示位置或分块因子。两者的秩一致时，形状运算会把每个因子应用到对应维度。

整数和普通元组都必须初始化，之后不能重新赋值。这项限制针对用于控制和形状计算的名字。缓冲区元素则不同：加法示例会写入 \co@output.at(i, j)@，矩阵乘法也会更新累加器。区分这两种角色，就不会把不可变索引误当成可变累加器。



\section{带形状数据与存储说明符(storage specifier)}
\label{sec:2-6}

\subsection*{完整类型}

缓冲区声明需要同时给出元素类型和形状。形状确定缓冲区包含多少个位置，基础类型确定每个位置存储什么。下面的例子为半精度数据复用命名形状，为单精度数据直接写出形状。

\begin{lstlisting}[language=croqtile,caption={完整类型}]
ndims : [20, 15];
f32 [10, 10] d0;
f16 [ndims] d1;
\end{lstlisting}

基本元素类型包括：

\begin{itemize}[leftmargin=1.6em]
\item 无符号整数：\co@u8@ / \co@u16@ / \co@u32@
\item 有符号整数：\co@s8@ / \co@s16@ / \co@s32@
\item 浮点数：\co@f16@ / \co@bf16@ / \co@f32@
\end{itemize}

\co@s32@ 是带形状数据的元素类型；用于程序控制的 \co@int@ 是独立的标量类型，两者用途不同。


\subsection*{存储说明符}

分块可以由协作线程共同使用，也可以由单个线程独占。元素类型和形状无法表达这种区别，需要由存储说明符补充。

这些示例使用以下三种存储说明符：

\begin{itemize}[leftmargin=1.6em]
\item \co@global@
\item \co@shared@
\item \co@local@
\end{itemize}

存储说明符写在元素类型之前：

\begin{lstlisting}[language=croqtile,caption={存储说明符}]
ndims : [20, 15];
local f32 [10, 10] d0;
shared f16 [ndims] d1;
\end{lstlisting}

应根据使用缓冲区的代码选择存储说明符。默认位置取决于执行上下文和目标平台；若搬运操作或协作组需要特定的存储位置，应显式声明。


声明缓冲区时，需要给出元素类型与形状；需要指定存储位置时，再加上存储说明符。


\subsection*{定义带形状数据}

下面声明一个带形状的缓冲区：

\begin{lstlisting}[language=croqtile,caption={带形状数据的声明}]
f32 [32, 16] data;
\end{lstlisting}

它的形状为 \co@[32, 16]@，元素类型为 \co@f32@。形状也可以先单独命名：

\begin{lstlisting}[language=croqtile,caption={独立的形状}]
shape : [72, 14];
s32 [shape] data0;              // using the named mdspan `shape`
u32 [shape(0), shape(1)] data1; // same shape as the above
u16 [shape + 6] data2;          // shaped as [78, 20]
s16 [shape, 8] data3;           // shaped as [72, 14, 8]
\end{lstlisting}

\co@data0@ 使用名为 \co@shape@ 的形状，并以 32 位整数为元素。\co@[]@ 内还可以使用形状运算的结果。

元素类型也称基础类型，常见类型如下：

\begin{itemize}[leftmargin=1.6em]
\item 无符号整数：\co@u8@、\co@u16@、\co@u32@
\item 有符号整数：\co@s8@、\co@s16@、\co@s32@
\item 浮点数：\co@f16@、\co@bf16@、\co@f32@
\end{itemize}

\co@u@ 和 \co@s@ 分别表示无符号与有符号整数，数字表示位宽。\co@f16@ 是半精度浮点，\co@bf16@ 是脑浮点，常用于机器学习计算。

\co@s32@ 用于声明数据元素，\co@int@ 用于程序控制。虽然位宽相同，它们的声明方式不同：

\begin{lstlisting}[language=croqtile,caption={被拒绝的基础类型}]
s32 a;    // a control integer, the same as 'int a'
s32 [4] b; // a buffer of four elements
\end{lstlisting}

两种写法都能通过编译：带形状的声明得到缓冲区，不带形状的声明得到普通控制整数；区别在形状，不在类型名。


要明确缓冲区的存储位置，可以加上说明符：

\begin{lstlisting}[language=croqtile,caption={存储说明符}]
global f32 [32, 7, 2] a;
shared u8 [512, 144] b;
local u8 [72, 1024] c;
\end{lstlisting}

存储位置由编译模式和目标平台约定决定。在这里的主机与设备协同场景中，未标注的缓冲区使用主机存储；设备侧存储位置则显式指定。

\begin{itemize}[leftmargin=1.6em]
\item \co@global@：指设备的全局存储。
\item \co@shared@：指设备块的共享存储。
\item \co@local@：指线程私有存储。
\end{itemize}

不同平台的含义和声明约束可能不同，应结合当前编译环境使用。


\subsection*{初始化}

缓冲区可以在声明时初始化：

\begin{lstlisting}[language=croqtile,caption={初始化}]
local s32 [17, 6] b1 {0};       // elements are initialized to 0
shared f32 [128, 16] b2 {3.14f}; // elements are initialized to 3.14f
\end{lstlisting}

变量名后的括号给出初始值。本例的初始化方式把所有元素设为同一个值，适合准备清零的输出或临时缓冲区。读取声明时，可以依次看存储说明符、元素类型、形状，再看名称后的初始化表达式。


\subsection*{声明参数与返回值}

带形状数据可以作为函数参数或返回值：

\begin{lstlisting}[language=croqtile,caption={显式的返回形状}]
__co__ f16 [7, 8] foo(f32 [16, 17, 5] input) {...}
\end{lstlisting}

参数与返回值的声明类似数据声明，但这里不能附带存储说明符或初始化表达式，也不能对已有 \co@mdspan@ 直接做运算。

参数的 \co@.span@ 返回它的形状，可以用来声明同形状的缓冲区：

\begin{lstlisting}[language=croqtile,caption={推导出的返回形状}]
__co__ auto foo(f32 [16, 17, 5] input) {
  f32 [input.span / {4, 1, 5} ] buffer;  // declare a buffer with tiled shape
  // ...
}
\end{lstlisting}


\subsection*{数据声明小结}

声明确定元素类型、形状与存储位置，初始化表达式给出初始内容。掌握这些写法后，就可以为并行工作准备相应的缓冲区。


\subsection*{\texttt{global}：全局存储}

声明中的 \co@global@ 把缓冲区放入目标平台的全局存储，适合需要超出单个线程或协作块访问范围的数据。选择存储位置并不等于初始化内容，也不会自动把已有数据搬入其中：

\begin{lstlisting}[language=croqtile,caption={全局存储声明}]
global f32 [100, 200] matrix;
\end{lstlisting}


\subsection*{\texttt{shared}：共享存储}

\co@shared@ 缓冲区为块内线程提供共同的暂存区域。在分块矩阵乘法中，一个块可以装入一次操作数分块，再让多个线程读取。读取必须在装入完成后开始，下次装入也必须等到这些读取结束。

\begin{lstlisting}[language=croqtile,caption={共享存储声明}]
shared f32 [10, 10] tile_data;
\end{lstlisting}


\subsection*{\texttt{local}：私有存储}

\co@local@ 表示线程私有存储，适合不需要在线程间共享的小数组或临时数据：

\begin{lstlisting}[language=croqtile,caption={私有存储声明}]
local f16 [5] thread_local_data;
\end{lstlisting}


\subsection*{默认行为与用法}

省略说明符时，存储位置由目标平台和所在执行上下文决定。入门代码这样写很方便，但仅凭一个没有说明符的名字，无法判断哪些线程能共享其内容。追踪数据流时，应把声明与包围它的并行区域一起阅读。

\begin{lstlisting}[language=croqtile,caption={不带说明符的声明}]
f32 [100, 100] d0;
\end{lstlisting}

选择 \co@shared@ 或 \co@local@ 时，将说明符写在元素类型之前：

\begin{lstlisting}[language=croqtile,caption={显式选择 local 存储}]
local f32 [10, 10] d0;     // Small data, local to the thread
shared f16 [ndims] d1;     // Intermediate data, shared between threads in a block
\end{lstlisting}


\subsection*{存储位置示例}

下面把三种存储位置放在同一个示例中：

\begin{lstlisting}[language=croqtile,caption={示例代码：使用存储说明符}]
ndims : [20, 15];
local f32 [10, 10] d0;
shared f16 [ndims] d1;
global f32 [100, 200] matrix;
// Perform computations on these arrays, utilizing the appropriate storage space
dma.copy d0 => d1;
dma.copy d1 => matrix;
\end{lstlisting}

这些声明展示三种存储位置，复制语句示意它们之间的数据流向。要把这段示意写成实际搬运，需要选择形状兼容的源视图和目标视图，因为这里声明的完整数组大小不同。存储说明符只选择位置，并不会使不同形状的缓冲区可以相互替代。



\section{迭代：\texttt{parallel-by}、\texttt{with-in} 与 \texttt{foreach}}
\label{sec:2-7}

\subsection*{Croqtile 的两层并行模型}

阅读并行代码时，先区分两个问题：哪些工作彼此独立，以及这些工作映射到哪里。

\begin{itemize}[leftmargin=1.6em]
\item \co@parallel@：声明彼此独立、可以并行执行的工作。
\item 空间说明符，如 \co@: block@、\co@: thread@、\co@: group@：指定并行轴的执行位置。
\end{itemize}

\co@foreach@ 表示顺序迭代，适合累加、归约或有前后依赖的处理；\co@parallel@ 表示独立的工作实例。

\begin{lstlisting}[language=croqtile,caption={Croqtile 的两层并行模型}]
parallel {px, py} by [8, 16] : block
  foreach {tile_k} in [16]               // sequential -- K-dimension loop
    parallel {qx, qy} by [16, 16] : thread  // concurrent -- mapped to threads
      foreach k in [16]                      // sequential -- inner accumulation
        output += lhs * rhs;
\end{lstlisting}

图~\ref{fig:f05-fig2_logical_vs_physical} 展示了空间说明符如何把各个轴映射到执行资源。

\begin{figure}[htbp]\centering
\includegraphics[width=\linewidth,height=0.56\textheight,keepaspectratio]{f05-fig2_logical_vs_physical_light.png}
\caption{空间说明符将 parallel 轴映射到 block 和 thread；其中的 foreach 循环仍按顺序执行。}
\label{fig:f05-fig2_logical_vs_physical}
\end{figure}

未写空间说明符时，由编译器选择默认映射；需要明确控制时，再标注 \co@: block@、\co@: thread@ 等说明符。

\co@parallel@ 和 \co@foreach@ 都可以使用多维索引。确定外层与内层范围后，组合运算符 \co@#@ 将它们的坐标连接成完整张量中的位置。范围参与这项计算，因此修改迭代划分时，分块选择也会随之改变。


\subsection*{\texttt{: block}——并行块}

\begin{lstlisting}[language=croqtile,caption={: block——并行块}]
parallel {px, py} by [8, 16] : block
\end{lstlisting}

这里声明 8 × 16 的网格，共 128 个块。每个块都应能独立完成分配给自己的输出分块；这个示例没有提供块间通信或同步协议。

本例输出形状为 \co@[128, 256]@，每块处理 \co@[16, 16]@，因此网格大小为 128/16 × 256/16 = 8 × 16。


\subsection*{\texttt{: thread}——块内线程}

\begin{lstlisting}[language=croqtile,caption={: thread——块内线程}]
parallel {qx, qy} by [16, 16] : thread
\end{lstlisting}

这一层在每个外层块内组织 16 × 16，共 256 个线程。

分块内的逐元素工作可以在这一层展开。


\subsection*{嵌套说明符}

块级和线程级工作可以嵌套声明：

\begin{lstlisting}[language=croqtile,caption={嵌套说明符}]
parallel {px, py} by [8, 16] : block {
  parallel {qx, qy} by [16, 16] : thread {
    // 128 blocks x 256 threads = 32,768 threads total
  }
}
\end{lstlisting}

外层划分块网格，内层划分块内线程。\co@px, py@ 表示同一并行层次中的两个坐标轴。


\subsection*{\texttt{with tile\_k in 16}}

语句 \co@with tile_k in 16@ 为 \co@tile_k@ 绑定 0 到 15 的范围，上界为 \co@#tile_k = 16@，但它本身不会把语句体执行 16 次。遍历由后续的 \co@foreach tile_k@ 完成。将绑定与循环分开后，流水线就可以先在循环外选择第 0 块，再遍历其余分块。


\subsection*{\texttt{foreach tile\_k(1:)}：切片迭代}

范围表达式 \co@(1:)@ 从索引 1 开始，保留原有的排他上界。共有 16 块时，它遍历 1 到 15。若前序阶段已装入第 0 块，循环就只需发起其余分块的装入，同时计算上一块；最后装入的一块留到循环结束后处理。


\subsection*{并行区域：\texttt{parallel-by} 块}

\co@parallel-by@ 用同一段代码描述多个独立工作实例。每个实例通过不同的索引处理自己的数据。

基本形式是在 \co@parallel ... by ...@ 后写出代码块：

\begin{lstlisting}[language=croqtile,caption={并行区域：parallel-by 块}]
parallel p by 6 {
  // SPMD code
}
\end{lstlisting}

本例包含 6 个实例，\co@p@ 分别取 \co@0@ 到 \co@5@。可以把它理解为 6 次彼此独立的迭代，执行先后顺序不固定。

\co@p@ 的取值范围为 \co@[0, 6)@，因此称为有界整数。分块选择等操作会同时使用它的当前值和范围。


\subsection*{\texttt{with-in} 块与 \texttt{where} 子句}

\co@with-in@ 将一个索引元组与形状范围绑定：

\begin{lstlisting}[language=croqtile,caption={索引 with-in 块}]
with index in [10, 10] {
  // index is ituple with 2 elements
}
\end{lstlisting}

\co@index@ 包含两个元素，也可以分别给元素命名：

\begin{lstlisting}[language=croqtile,caption={索引元组}]
with {x, y} in [10, 10] {...}
\end{lstlisting}

元组本身和各个元素可以同时命名：

\begin{lstlisting}[language=croqtile,caption={由元组导出的索引}]
with index = {x, y} in [10, 10] {...}
\end{lstlisting}

这里的 \co@index@ 是有界整数元组。

\co@with-in@ 只声明有界索引，不启动并行工作。需要遍历这些索引时，再使用 \co@foreach@。

\co@where@ 可以声明索引之间的对应关系：

\begin{lstlisting}[language=croqtile,caption={嵌套的有界组}]
with {m, n} in [M, N], {n_p, k} in [N_P, K] where n_p <-> n {
  // matmul implements with m,n,K. n_p is no long useful.
}
\end{lstlisting}

\co@n <-> n_p@ 约定两个索引在迭代中对应同一取值，因此可以在需要该索引的位置使用相应名字。


\subsection*{\texttt{foreach} 块}

有界整数或元组定义后，可以用 \co@foreach@ 遍历其范围。这里把范围声明与循环体分开：\co@with-in@ 说明可用的位置，\co@foreach@ 再按顺序逐个访问。循环体中的索引代表当前正在处理的位置：

\begin{lstlisting}[language=croqtile,caption={foreach 块}]
with x in [10] {
  foreach x {
    // do something with each x
  }
}
\end{lstlisting}

每次迭代在循环体内为有界变量提供当前值，上界则保持不变。选择数据时使用当前值，推导分块数量时使用上界。若改为声明并行索引，则是给每个独立实例分配一个值，而不是按顺序遍历这些值。


\subsection*{具名与不具名的并行索引}

下面用具名索引声明并行工作：

\begin{lstlisting}[language=croqtile,caption={具名的并行索引}]
parallel p by 6 {
  // SPMD (Single Program, Multiple Data) code
}
\end{lstlisting}

花括号中的代码由多个实例执行。本例的组成部分如下：

\begin{itemize}[leftmargin=1.6em]
\item \co@parallel@：声明独立的工作实例。
\item \co@p@：标识当前实例的有界索引。
\item \co@by 6@：声明 6 个实例，索引范围为 \co@0@ 到 \co@5@。
\end{itemize}

如果代码不需要引用实例索引，可以省略变量名：

\begin{lstlisting}[language=croqtile,caption={不具名的并行索引}]
parallel by 2 { ... }
\end{lstlisting}

这里仍有两个工作实例，但代码不能用省略的索引来区分数据位置。


\subsection*{多层并行}

可以嵌套声明不同的并行层次：

\begin{lstlisting}[language=croqtile,caption={嵌套的两级并行}]
parallel p by 2 {
  // parallel-level-0
  parallel q by 12 {
    // parallel-level-1
  }
}
\end{lstlisting}

本例的两层范围分别为 \co@2@ 和 \co@12@，共有 24 种索引组合。\co@p@ 属于外层，\co@q@ 属于内层，每个组合标识一份工作。

层次划分也决定代码在哪个范围内执行，以及哪些工作需要协调。组织代码时，应让协作关系与并行层次相匹配。

两层之间没有其他语句时，也可以合写为一条声明：

\begin{lstlisting}[language=croqtile,caption={合写的两级并行}]
parallel p by 2, q by 12 { ... }
\end{lstlisting}

逗号分隔的写法与前面的嵌套声明表达相同的两级并行。


\subsection*{同一并行层次的多维坐标}

同一并行层次可以用多个坐标轴描述；支持的轴数由目标平台决定：

\begin{lstlisting}[language=croqtile,caption={子级并行性}]
parallel {px, py, pz} by [1, 1, 2] {
  // higher synchronization cost
  // ...
  parallel t_index = {qx, qy} by [3, 4] {
    // lower synchronization cost
    // ...
  }
  // ...
}
\end{lstlisting}

本例外层使用 \co@px, py, pz@ 三个轴，内层使用 \co@qx, qy@ 两个轴。内层索引还可以整体命名为 \co@t_index@，各轴的范围由对应的形状给出。

元组中第一个元素对应最低有效的并行轴。使用多维索引时，需要保持这一顺序与数据坐标的对应关系。


\co@with-in@ 定义索引范围，\co@foreach@ 按该范围依次迭代。两者一起描述循环的边界和循环体。


\subsection*{用有界索引描述循环}

C++ 的循环由初值、条件和步进控制。Croqtile 的循环则围绕有界索引组织。

\co@foreach@ 的基本形式如下：

\begin{lstlisting}[language=croqtile,caption={一个 foreach 循环}]
foreach index {
  // loop body
}
\end{lstlisting}

如果 \co@index@ 的上界为 \co@6@，它对应下面的 C++ 循环：

\begin{lstlisting}[language=croqtile,caption={等价的 C++ 循环}]
for (int index = 0; index < 6; ++index) {
  // loop body
}
\end{lstlisting}

有界索引把循环范围明确写在程序中，便于分块选择与数据搬运使用。

\co@parallel-by@ 中的索引标识并行实例，不可修改；它可以作为 \co@foreach@ 的上界，但不能作为被遍历的变量。要声明循环索引，应使用 \co@with-in@。





\subsection*{定义有界变量}

\co@with-in@ 可以定义有界整数或有界整数元组，不创建额外的执行实例：

\begin{lstlisting}[language=croqtile,caption={一个有界变量}]
with x in 128 {
  // 'x' is a bounded integer
}
with y in [512] {
  // 'y' is a bounded ituple
}
with index in [10, 10] {
  // 'index' is a bounded ituple
}
\end{lstlisting}

本例中，\co@x@ 的范围是 128；\co@y@ 是一维元组，范围为 \co@[512]@；\co@index@ 是二维元组，范围为 \co@[10, 10]@。元组的元素也可以分别命名。

\begin{lstlisting}[language=croqtile,caption={一组有界变量}]
with {x, y} in [10, 10] {
  // x and y are explicitly named elements of the ituple
}
with index = {x, y} in [10, 10] {
  // 'x', 'y', and 'index' can be used within the block
}
\end{lstlisting}

多条声明可以用逗号合写：

\begin{lstlisting}[language=croqtile,caption={由元组导出的变量}]
with index = {x, y} in [10, 10], idx in [100, 10] { }
\end{lstlisting}

\co@index@ 与 \co@idx@ 只在随后的 \co@with-in@ 作用域内有效。





用 \co@foreach@ 遍历已经定义的有界变量：

\begin{lstlisting}[language=croqtile,caption={索引一个形状}]
with index in [6] {
  foreach index {
    // Perform operations with each index
  }
}
\end{lstlisting}

本例迭代 6 次，索引从 \co@0@ 递增到 \co@5@。

也可以遍历多维元组：

\begin{lstlisting}[language=croqtile,caption={索引元组}]
with index = {x, y} in [6, 17] {
  foreach index { }
}
\end{lstlisting}

对 \co@index@ 的遍历对应下面的嵌套循环：

\begin{lstlisting}[language=croqtile,caption={生成的 C++ 循环}]
for (int x = 0; x < 6; ++x)
  for (int y = 0; y < 17; ++y) { }
\end{lstlisting}

元组中的索引从左到右，对应循环从外到内的顺序。

\co@foreach@ 后列出多个有界变量时，采用相同的嵌套规则：

\begin{lstlisting}[language=croqtile,caption={若干有界组}]
with index = {x, y} in [6, 17], iv in [128] {
  foreach iv, index { }
}
\end{lstlisting}

下面的等价代码把 \co@iv@ 放在最外层，再依次遍历 \co@index@ 的两个维度。

\begin{lstlisting}[language=croqtile,caption={生成的内层循环}]
for (int iv = 0; iv < 128; ++iv)
  for (int x = 0; x < 6; ++x)
    for (int y = 0; y < 17; ++y) { }
\end{lstlisting}


\subsection*{直接声明循环范围}

\begin{lstlisting}[language=croqtile,caption={一个 foreach 循环}]
foreach x in 128 { }
foreach idx in [10, 20] { }
foreach {y, z} in [8, 16] { }
\end{lstlisting}

如果一个索引只用于紧接着的循环，可以直接在 \co@foreach@ 后给出范围，省去单独的 \co@with-in@ 声明。两种写法的区别在于声明位置，循环范围与遍历顺序不变：

\begin{lstlisting}[language=croqtile,caption={它展开成的 with-in 形式}]
with x in 128 {
  foreach x { }
}
with idx in [10, 20] {
  foreach idx { }
}
with {y, z} in [8, 16] {
  foreach y, z { }
}
\end{lstlisting}


\subsection*{由有界整数导出循环}

需要跳过起始元素、提前结束或按步长迭代时，可以在有界整数后附加范围表达式：

\begin{lstlisting}[language=croqtile,caption={一组有界变量}]
with {x, y} in [6, 17] {
  foreach x, y(1::) { }
}
\end{lstlisting}

\co@y(1::)@ 对应下面的循环：

\begin{lstlisting}[language=croqtile,caption={导出的 C++ 循环}]
for (int x = 0; x < 6; ++x)
  for (int y = 1; y < 17; ++y) { }
\end{lstlisting}

循环从 \co@1@ 开始。范围表达式的三个字段分别表示下偏移、上偏移和步长：

\begin{lstlisting}[language=croqtile,caption={有界变量的语法}]
  bounded-variable(lower-offset:upper-offset:stride)
\end{lstlisting}

它保留原有的索引范围，并调整实际遍历方式：

\begin{itemize}[leftmargin=1.6em]
\item 初值为索引下界加下偏移；有界整数的下界为 \co@0@。
\item 结束位置为索引上界加上偏移，不包含该位置。
\item 每次迭代按给定步长递增。
\end{itemize}

例如 \co@y(1:-1:2)@ 对应 \co@for (y = 0 + 1; y < 17 - 1; y += 2)@。省略字段时，下偏移和上偏移默认为 \co@0@，步长默认为 \co@1@。

范围表达式适用于有界整数，不能直接用于有界整数元组。


\subsection*{循环与范围小结}

有界索引给出循环范围，范围表达式调整实际遍历的起点、终点与步长。

跳过已预先处理的分块等写法，会在 3.6.2 节的多缓冲模式中使用。



\section{分块选择：\texttt{chunkat}、\texttt{subspan} 与 \texttt{view}/\texttt{from}}
\label{sec:2-8}

\subsection*{\texttt{chunkat(...)}：等分}

\begin{lstlisting}[language=croqtile,caption={chunkat(...)：等分}]
lhs.chunkat(tr, tc)
\end{lstlisting}

表达式 \co@chunkat@ 根据对应索引的上界划分各维。输入为 \co@[64, 128]@、并行声明为 \co@parallel {tr, tc} by [4, 8]@ 时，选中的分块大小是 \co@[16, 16]@。\co@tr@ 和 \co@tc@ 的当前值决定选哪块，上界决定块有多大。

已有并行或循环索引明确了分块方式时，可以用 \co@chunkat@ 选取当前分块。本例中，\co@lhs@ 的形状为 \co@[64, 128]@，\co@tr@、\co@tc@ 的范围分别为 \co@4@、\co@8@，所以每块为 \co@[16, 16]@。索引当前值再决定选中哪一块。

\begin{figure}[htbp]\centering
\includegraphics[width=\linewidth,height=0.56\textheight,keepaspectratio]{f04-fig3_chunkat_light.png}
\caption{对按 [4, 8] 网格划分的 [64, 128] 张量，chunkat(1, 3) 选中第 16–31 行、第 48–63 列。}
\end{figure}


\subsection*{\texttt{subspan(...).at(...)}：指定分块形状与坐标}

\begin{lstlisting}[language=croqtile,caption={subspan(...).at(...)：分块网格上显式指定分块范围}]
lhs.subspan(64, 64).at(block_m, iv_k)
\end{lstlisting}

表达式 \co@subspan@ 显式给出分块尺寸，\co@.at(...)@ 给出分块坐标。对于 \co@[16, 16]@ 分块，坐标 \co@(2, 3)@ 的元素起点是 \co@(32, 48)@。分块大小由算法固定选择，而不是从并行实例数推导时，适合使用这种形式。


\subsection*{\texttt{view(...).from(...)}：指定区域形状与元素偏移}

\begin{lstlisting}[language=croqtile,caption={view(...).from(...)：显式分块范围与显式元素偏移}]
lhs.view(TILE_M, TILE_K).from(offset_m, offset_k)
\end{lstlisting}

形式 \co@view(...).from(...)@ 为窗口给出显式形状和按元素计量的起点。从 \co@(37, 50)@ 开始的 \co@[16, 16]@ 窗口保留这个起点，不会将其取整到最近的分块边界。

起始位置以元素偏移给出，或访问区域不适合规则等分时，可以使用 \co@view(...).from(...)@。\co@view@ 确定窗口大小，\co@from@ 确定它从哪个元素开始；窗口位置因此不必落在分块网格的整数坐标上。

三种写法分别适用于：

\begin{itemize}[leftmargin=1.6em]
\item \co@chunkat(...)@：由有界索引确定等分方式和分块坐标。
\item \co@subspan(...).at(...)@：显式给出分块形状与分块坐标。
\item \co@view(...).from(...)@：显式给出区域形状与起始元素偏移。
\end{itemize}


\subsection*{组合：\texttt{outer \# inner}}

\begin{lstlisting}[language=croqtile,caption={组合：outer \# inner}]
output.at(tr # i, tc # j)
\end{lstlisting}

中缀 \co@#@ 将外层分块坐标和块内坐标组合成原张量中的坐标。

对于外层分块索引 \co@tr@ 和块内索引 \co@i@，\co@tr # i@ 的值为 \co@tr * #i + i@。每块有 16 行时，第 2 块的块内第 5 行对应全局第 37 行。组合后的上界是两个上界的乘积。内层上界决定步幅，因此操作数顺序不能交换。

图~\ref{fig:f04-fig_compose} 展示了分块坐标与块内坐标如何合成为全局坐标。

\begin{figure}[htbp]\centering
\includegraphics[width=\linewidth,height=0.56\textheight,keepaspectratio]{f04-fig_compose_light.png}
\caption{每块含 16 个元素时，tile = 2 与块内偏移 i = 3 组合得到全局位置 2 × 16 + 3 = 35。}
\label{fig:f04-fig_compose}
\end{figure}


\subsection*{范围：\texttt{\#name}}

\begin{lstlisting}[language=croqtile,caption={范围：\#name}]
foreach i in [128 / #tile]
\end{lstlisting}

前缀 \co@#tile@ 返回 \co@tile@ 的范围。声明为 \co@parallel tile by 8@ 时，\co@#tile@ 就是 \co@8@。

图~\ref{fig:f04-fig_extent} 展示了 \texttt{\#} 前缀的两种用法。

\begin{figure}[htbp]\centering
\includegraphics[width=\linewidth,height=0.56\textheight,keepaspectratio]{f04-fig_extent_light.png}
\caption{前缀 \#tile 返回轴的范围大小，此处为 8；中缀 tile \# i 将分块坐标与块内偏移组合。}
\label{fig:f04-fig_extent}
\end{figure}


\subsection*{用有界索引选择分块}

\co@chunkat@ 表达式从已有带形状数据中选择视图。每个有界参数提供两项信息：上界决定分块数量，当前值选择其中一块。选择视图本身不搬运元素，实际搬运由外层 DMA 语句执行。

\begin{lstlisting}[language=croqtile,caption={有界整数元组／整数与 chunkat 运算}]
global f32 [6, 10, 100] data;
parallel p by 6 {
  with index = {x, y} in [10, 10] {
    // for every data move, the stride into 'data' is
    //    stride = p*1000 + x * 100 + y * 10
    //
    // for each chunk, the dimensioned size for the movement is {1, 1, 10}
    f = dma.copy data.chunkat(p, index) => local;
  }
}
\end{lstlisting}

本例输入形状为 \co@[6, 10, 100]@，索引范围把它划为 6 × 10 × 10 个分块。每块在最后一维包含 10 个连续元素，并由搬运语句送入局部存储。


在 DMA 语句中，\co@chunkat@ 直接描述要搬运的分块，减少手工计算地址偏移的工作。


\subsection*{上界运算}

有界变量同时包含当前索引和范围。循环使用索引，分块与缓冲区声明则常常需要它的上界。

在有界变量前加 \co@#@ 可以取出上界：

\begin{lstlisting}[language=croqtile,caption={上界运算}]
with size in [32] {
  local u8 [#size, 16] a; // buffer shape: [32, 16]
  local u8 [size, 16] b; // error in shape: [0, 16]
}
\end{lstlisting}

\co@#size@ 给出范围 \co@[32, 16]@，可以用于声明缓冲区 \co@a@。\co@size@ 本身是当前索引，不能在这里替代范围。


\subsection*{用 \texttt{chunkat} 表达式对数据分块}

搬运可以使用整个缓冲区，但分块算法通常只需要当前实例负责的区域。把 \co@chunkat@ 放在源表达式中，就能直接描述这片区域：

\begin{lstlisting}[language=croqtile,caption={用 chunkat 表达式对数据分块}]
__co__ void foo(f32 [36] input) {
  parallel p by 6 : block {
    dma.copy input.chunkat(p) => shared;
  }
}
\end{lstlisting}

\co@.chunkat(...)@ 根据有界索引选取当前分块，可以理解为先按范围分块，再按当前索引选取。

\begin{figure}[htbp]\centering
\includegraphics[width=\linewidth,height=0.56\textheight,keepaspectratio]{f29-chunkat.png}
\caption{chunkat 表达式各部分的语义}
\end{figure}

\co@input.chunkat(p)@ 对应 \co@input.chunk(#p).at(p)@，各部分的作用如下：

\begin{itemize}[leftmargin=1.6em]
\item \co@input@：要分块的数据。
\item \co@#p@：分块因子。本例为 \co@6@，将输入分成 6 块。
\item \co@.chunk(#p)@：确定每块形状，本例为 \co@[6]@。
\item \co@.at(p)@：选择当前块。例如 \co@p == 1@ 时选择元素 \co@[6, 12)@，\co@p == 5@ 时选择 \co@[30, 36)@。
\end{itemize}

\co@chunkat@ 把分块因子与分块索引放在同一表达式中，避免分别维护两组信息。


\subsection*{使用多个有界变量}

多个有界变量可以一起指定分块：

\begin{lstlisting}[language=croqtile,caption={使用多个有界变量}]
__co__ void foo(f32 [36, 14, 8] input) {
  parallel p by 6 : block {
    with index in [7, 4] {
      foreach index {
        dma.copy input.chunkat(p, index) => shared;
      }
    }
  }
}
\end{lstlisting}

本例将 \co@p@ 与二维元组 \co@index@ 组合起来。各维范围为 \co@6@、\co@7@、\co@4@，所以 \co@[36, 14, 8]@ 的输入被划为 \co@[6, 2, 2]@ 的小块。

外层 6 个并行实例各处理一个 \co@[6, 14, 8]@ 区域，内部再迭代 7 × 4 次，每次把一个 \co@[6, 2, 2]@ 分块搬入 \co@shared@。

同一组范围现在同时描述迭代空间与所选数据。划分改变时，编译器从这些范围推导新的分块尺寸。检查搬运时，可以先将输入形状除以上界，得到分块大小，再把分块坐标乘以相应尺寸，还原元素起点。


\subsection*{分块选择小结}

\co@chunkat@ 用有界索引选取分块，既可用于 DMA 的源，也可用于目的区域。

3.6.1 节将分块、迭代与搬运组合成完整示例。



\section{数据搬运：DMA 语句与期值(future)}
\label{sec:2-9}

\subsection*{DMA 语句族}

分块搬运使用统一的基本形式：

\begin{lstlisting}[language=croqtile,caption={DMA 的一般形式}]
future = engine.operation.modifiers source => destination;
\end{lstlisting}

各部分含义如下：

\begin{table}[htbp]\centering\small
\begin{tabularx}{\linewidth}{@{}lX@{}}
\toprule
组成部分 & 含义 \\
\midrule
\co@future@ & 可选的传输结果句柄。用 \co@future.data@ 访问已搬运分块，用 \co@future.span@ 查询其形状。 \\
\co@operation@ & \co@copy@、\co@transp<...>@ 或 \co@pad<...>@。 \\
\co@modifiers@ & 可选标志，如 \co@.async@、\co@.zfill@ 或 \co@.swiz<N>@。 \\
\co@source => destination@ & 源区域与目的存储位置或区域。 \\
\bottomrule
\end{tabularx}
\end{table}

最简单的拷贝示例如下：

\begin{lstlisting}[language=croqtile,caption={一个具体的拷贝}]
f = dma.copy input.chunkat(blk) => shared;
g = dma.transp<1, 0> matrix => local;
h = dma.pad<{2, 1}, {3, 2}, {0, 0}, 0> tile => shared;
\end{lstlisting}

搬运对象是带形状的数据区域，可以一次描述整个分块的传输。\co@=>@ 左侧写源区域，右侧写目的区域或存储位置；目的只写 \co@shared@ 时，编译器根据源形状准备缓冲区，后续代码可以通过期值引用它。


\subsection*{示例：先搬运分块，再计算加法}

\begin{lstlisting}[language=croqtile,caption={第一个示例：经由共享存储的分块加法}]
__co__ s32 [64, 128] tiled_add_2d(s32 [64, 128] lhs, s32 [64, 128] rhs) {
  s32 [lhs.span] output;
  parallel {tr, tc} by [4, 8] : block {
    lhs_load = dma.copy lhs.chunkat(tr, tc) => shared;
    rhs_load = dma.copy rhs.chunkat(tr, tc) => shared;
    foreach {i, j} in lhs_load.span
      output.at(tr # i, tc # j) =
          lhs_load.data.at(i, j) + rhs_load.data.at(i, j);
  }
  return output;
}
\end{lstlisting}

下面仍计算逐元素加法，先把输入分块搬入共享存储：

\begin{itemize}[leftmargin=1.6em]
\item \co@chunkat(tr, tc)@ 选择输入分块。
\item \co@dma.copy ... => shared@ 将分块搬入共享缓冲区。
\item \co@tr # i@ 与 \co@tc # j@ 将块内坐标组合成输出坐标。
\end{itemize}

程序依次选择分块、搬入缓冲区，再对其中的元素计算。两个输入的同一分块分别由 \co@lhs_load@ 和 \co@rhs_load@ 引用；内部循环使用搬运后的形状遍历元素，最后把块内索引组合成原输出中的位置。

\begin{figure}[htbp]\centering
\includegraphics[width=\linewidth,height=0.56\textheight,keepaspectratio]{f04-fig2_tiled_add_light.png}
\caption{输入形状为 [64, 128]，分块网格为 [4, 8]。分块 (1, 3) 先装入两个 [16, 16] 操作数，再逐元素相加。}
\end{figure}


\subsection*{期值、\texttt{.data} 与 \texttt{.span}}

DMA 传输可以返回一个期值句柄：

\begin{lstlisting}[language=croqtile,caption={由一次拷贝得到的期值}]
f = dma.copy input => shared;
\end{lstlisting}

它提供两个成员：

\begin{itemize}[leftmargin=1.6em]
\item \co@f.data@：传输的目的缓冲区。
\item \co@f.span@：目的缓冲区的形状。
\end{itemize}

使用方式如下：

\begin{lstlisting}[language=croqtile,caption={读取 data 与 span}]
f = dma.copy input => shared;
foreach idx in [f.span]
  output.at(idx) = f.data.at(idx);
\end{lstlisting}

即使拷贝是同步的，给结果命名也有用途：后续代码既能通过 \co@.data@ 访问目的缓冲区，又能通过 \co@.span@ 确定遍历范围。这样不必再单独维护一份缓冲区名称和形状。

图~\ref{fig:f04-fig_future_data} 展示了拷贝返回的句柄暴露了哪些信息。

\begin{figure}[htbp]\centering
\includegraphics[width=\linewidth,height=0.56\textheight,keepaspectratio]{f04-fig_future_data_light.png}
\caption{复制结果通过 .data 提供目标张量，通过 .span 提供其形状；异步复制必须完成后才能读取。}
\label{fig:f04-fig_future_data}
\end{figure}


\subsection*{同步拷贝与异步拷贝}

不加 \co@.async@ 时，拷贝完成后才执行后续语句：

\begin{lstlisting}[language=croqtile,caption={同步拷贝}]
f = dma.copy input => shared;
dma.copy f.data => output;
\end{lstlisting}

加上 \co@.async@ 后，可以在传输期间继续执行其他工作：

\begin{lstlisting}[language=croqtile,caption={异步拷贝}]
f = dma.copy.async input => shared;
wait f;
dma.copy f.data => output;
\end{lstlisting}

异步拷贝使用相同的形式：

\begin{lstlisting}[language=croqtile,caption={异步拷贝}]
f = dma.copy.async input => shared;
wait f;
dma.copy f.data => output;
\end{lstlisting}

期值可以在异步搬运尚未结束时就标识目标缓冲区。能够引用 \co@f.data@，并不意味着元素已经就绪。应在首次读取这些元素的操作之前执行 \co@wait f@，并在搬运过程中保持源和目标有效。

\co@.async@ 发起传输，\co@wait@ 确认完成；两者共同确定数据何时可以使用。


\subsection*{\texttt{dma.transp\textless{}...\textgreater{}}}

\begin{lstlisting}[language=croqtile,caption={dma.transp\textless{}...\textgreater{}}]
fa = dma.transp<1, 0, 2> a => local;
dma.copy fa.data => o;
\end{lstlisting}

\co@dma.transp@ 在搬运时置换维度，目的缓冲区使用变换后的布局。它把维度变换与数据搬运写在同一条语句里；读取代码时，既要看源和目的，也要看尖括号中约定的维度顺序。


\subsection*{\texttt{dma.pad\textless{}...\textgreater{}}}

\begin{lstlisting}[language=croqtile,caption={dma.pad\textless{}...\textgreater{}}]
f = dma.pad<{2, 1}, {3, 2}, {0, 0}, V>.async input => shared;
wait f;
dma.copy f.data => output;
\end{lstlisting}

\co@dma.pad@ 在搬运时填充目的区域。模板实参指定各维填充、内部间隔和填充值，适合需要显式边界值的情况。与普通拷贝相比，目的区域可能更大，后续计算应使用填充后得到的形状。


\subsection*{Tileflow 程序：编排数据搬运}

下面把迭代、分块选择和 DMA 放在同一个 Tileflow 程序中，观察数据从输入到输出的过程。

完整代码如下：

\begin{lstlisting}[language=croqtile,caption={tileflow 程序}]
__co__ s32 [6, 17, 128] ele_add(s32 [6, 17, 128] lhs, s32 [6, 17, 128] rhs) {
  s32 [lhs.span] output; // Use same shape as lhs
  parallel p by 6 {
    with index in [17, 4] { // Tiling factors
      foreach index {
        lhs_load = dma.copy lhs.chunkat(p, index) => local;
        rhs_load = dma.copy rhs.chunkat(p, index) => local;
        local s32 [lhs_load.span] l1_out;
        // Call kernel with loaded data
        call kernel(lhs_load.data, rhs_load.data, l1_out, |lhs_load.span|);
        // Store result back to output
        dma.copy l1_out => output.chunkat(p, index);
      }
    }
  }
  return output;
}
\end{lstlisting}

\co@__co__@ 函数接收形状为 \co@[6, 17, 128]@ 的两个 \co@s32@ 输入，并声明同形状的输出。

\co@parallel p by 6@ 声明 6 个独立实例，每个实例处理输入的一个外层区域。

\co@with-in@ 将 \co@index@ 绑定到范围 \co@[17, 4]@。它是二维有界整数元组，供内部循环使用。

\co@foreach index@ 对应下面的嵌套循环：

\begin{lstlisting}[language=croqtile,caption={等价的 C++ 循环}]
for (int x = 0; x < 17; x++)  // assume 'x' represents the 1st element of 'index'
  for (int y = 0; y < 4; y++) { ... }  // and 'y' represents the 2nd element of 'index'
\end{lstlisting}

循环中的 \co@dma.copy@ 负责搬运当前分块：

\begin{itemize}[leftmargin=1.6em]
\item \co@lhs_load@：传输结果句柄，用于引用目的缓冲区。
\item \co@dma.copy@：拷贝数据，\co@=>@ 左边是源，右边是目的。
\item \co@local@：让编译器分配线程私有的目的缓冲区。
\item \co@lhs.chunkat(p, index)@：根据范围 \co@6@、\co@17@、\co@4@ 选取 \co@[1, 1, 32]@ 分块，具体位置由索引当前值决定。例如索引 \co@(1, 16, 2)@ 选中对应的分块。
\end{itemize}

分块方式如下图所示。

\begin{figure}[htbp]\centering
\includegraphics[width=\linewidth,height=0.56\textheight,keepaspectratio]{f14-chunkat_3d_light.png}
\caption{形状 [6, 17, 128] 按范围 [6, 17, 4] 划分后，每块为 [1, 1, 32]。分块 (1, 16, 2) 从元素 (1, 16, 64) 开始。}
\end{figure}

\co@lhs@ 和 \co@rhs@ 的对应分块分别搬入设备的 \co@local@ 存储，供当前实例计算。

\co@local s32 [lhs_load.span] l1_out;@ 声明同分块形状的输出缓冲区。\co@call kernel@ 在其中完成计算，随后 DMA 将结果搬回主机侧输出区域。

每个实例迭代 17 × 4 次，每次处理一个 \co@[1, 1, 32]@ 分块。全部实例完成后，函数把输出交给主机调用方。


DMA 语句同时描述要搬运的数据区域和传输结果，结果可以用期值命名。


异步 DMA 允许程序在传输期间继续执行后续代码；使用传输结果前，需要等待完成。

基本写法如下：

\begin{lstlisting}[language=croqtile,caption={基本语法}]
future = dma-op src-expr => dst-expr;
\end{lstlisting}

\co@=@ 左边的期值命名传输结果。同步传输也可以命名，以便引用目的数据；不需要引用时可以省略。

\co@=@ 右边依次写搬运操作、源表达式和目的表达式。\co@=>@ 表示从源到目的的数据流向。


\subsection*{同步与异步}

异步复制是非阻塞的，搬运尚未完成时后续代码可以继续执行。同步复制则在语句执行结束前完成搬运。这些术语描述发起代码何时能够继续，并不表示可以安全读取未完成搬运的目标。

下面同时使用异步和同步拷贝：

\begin{lstlisting}[language=croqtile,caption={同步与异步}]
f0 = dma.copy.async data0 => shared; // non-blocking
f1 = dma.copy data1 => shared;       // blocking
...
wait f0;                             // synchronization
\end{lstlisting}

\co@f0@ 的拷贝在后台进行，程序继续执行；\co@f1@ 的拷贝是同步的，需要等它完成才能继续。等待 \co@f1@ 时，\co@f0@ 的传输可能仍在进行。

\co@wait f0;@ 等待对应传输完成。应在读取结果之前完成等待，并安排好所有异步传输的使用顺序。


\subsection*{数据表达式}

数据表达式可以引用已有缓冲区。目的也可以直接写存储位置，由编译器分配相应缓冲区：

\begin{lstlisting}[language=croqtile,caption={数据表达式}]
global f32 [32, 16] input;
shared f32 [32, 16] output;
dma.copy input => output;  // Copy to a user-declared buffer
dma.copy input => shared;  // Same, but to a compiler-allocated buffer
\end{lstlisting}

第二条语句以 \co@shared@ 为目的，编译器按源形状分配缓冲区，效果与第一条显式声明的写法相同。

源和目标位置也可以使用分块表达式。例如 \co@input.chunkat(p)@ 将复制限制在 \co@p@ 选中的区域，显式目标视图则限制写入位置。第 2.8 节说明这些视图怎样形成。


\subsection*{期值与等待}

期值命名传输结果。异步传输的期值还可以传给 \co@wait@，用来等待完成。

\co@wait@ 可以同时等待多个异步期值；同步传输已经完成，不能再用其期值等待。

\begin{lstlisting}[language=croqtile,caption={同步与异步期值的等待}]
f0 = dma.copy.async input0 => shared;
f1 = dma.copy.async input1 => shared;
f2 = dma.copy input2 => shared;
...
wait f0, f1;  // Multiple wait
wait f2;      // Error: cannot wait on a sync-dma
\end{lstlisting}

同步传输返回期值，是为了让后续代码取得目的缓冲区。主要成员如下：

\begin{itemize}[leftmargin=1.6em]
\item \co@.span@ 取出 DMA 目的缓冲区的 \co@mdspan@。
\item \co@.data@ 取出 DMA 目的缓冲区的引用（带形状数据）。
\end{itemize}

下面通过期值访问目的缓冲区：

\begin{lstlisting}[language=croqtile,caption={通过期值访问数据与形状}]
f0 = dma.copy.async input0 => shared;
f1 = dma.copy input2 => shared;
...
local f32 [f0.span] buffer;
call device_kernel(f1.data, |f1.span|);
\end{lstlisting}

\co@.span@ 返回目的形状，\co@.data@ 返回目的数据。示例中的元素计数运算还可以取得目的区域的总元素数。


\subsection*{数据搬运小结}

搬运操作确定传输方式，期值命名结果，\co@wait@ 明确异步数据何时就绪。读代码时，可以沿着这三个环节追踪数据流。

需要回顾分块区域的选择方式时，可查阅 2.8 节。



\section{协调：\texttt{inthreads}、事件与异步块}
\label{sec:2-10}

\subsection*{选择执行实例}

\co@inthreads@ 从当前并行作用域中选择执行实例。

下面只让索引不大于 2 的实例执行选定区域：

\begin{lstlisting}[language=croqtile,caption={tileflow 程序的并行视角}]
__co__  void foo() {
  // sequential code
  // start of the SPMD code region
  parallel p by 6 {
    // some uniform code for all threads
    // ...
    inthreads (p <= 2) {
      // divergent code for thread 0, 1, 2 (thread-mask: 000111)
    }
    // threads are forced to be synchronized
    // some uniform code for all threads
    // ...
  }
  // end of the SPMD code region
  // sequential code
}
\end{lstlisting}

可以把程序中的执行范围分成以下几类：

\begin{itemize}[leftmargin=1.6em]
\item 主机侧的顺序代码。
\item \co@parallel-by@ 内由多个实例执行的代码。
\item 并行实例共同执行的代码。
\item \co@inthreads@ 中仅由选定实例执行的代码。
\end{itemize}

本例的 \co@p <= 2@ 选择索引为 \co@0@、\co@1@、\co@2@ 的实例；其余实例跳过该区域。普通 \co@inthreads@ 用于按顺序执行的过滤区域，异步形式则使用 \co@.async@。

\co@inthreads@ 从已有并行区域中选择实例，使用位置和条件需要满足以下约束：

\begin{itemize}[leftmargin=1.6em]
\item 只能放在 \co@parallel-by@ 创建的并行作用域中。
\item 条件需要与并行索引有关。
\end{itemize}

编译器会检查条件表达式及其所在的并行作用域。


\subsection*{嵌套层次之间的执行范围}

嵌套并行层次之间的语句，不一定由内层的所有实例重复执行：

\begin{lstlisting}[language=croqtile,caption={隐式掩码}]
__co__  void foo() {
  parallel p by 6 {
    // implicit masking: 'inthreads (q == 0) { ... } '
    parallel q by 2 {
      // code here is for all threads
    }
    // implicit masking: 'inthreads (q == 0) { ... } '
  }
}
\end{lstlisting}

本例中，位于外层 \co@parallel p by 6@、内层 \co@parallel q by 2@ 之外的代码，每个外层实例只执行一份。可以理解为内层索引 \co@q == 0@ 的实例执行它；进入内层块后，才由该层的全部实例执行。


异步区域可以独立推进，但对另一区域的依赖仍需显式表达。事件标记这项依赖：一个区域完成所需工作后触发事件，另一区域在使用结果前等待。事件传递就绪信息，数据仍保存在这些区域共用的缓冲区中。


\subsection*{事件基础}

事件用于协调异步代码区域的执行顺序。例如，一个区域准备数据，另一个区域要在准备完成后才能读取。生产者触发事件，消费者等待事件，就把这项先后关系直接写进了程序。

事件有 \emph{SET}（已触发）和 \emph{UNSET}（未触发）两种状态，声明时为 \emph{UNSET}：

\begin{lstlisting}[language=croqtile,caption={声明一个共享事件}]
shared event e;
\end{lstlisting}

\co@shared@ 事件在并行作用域中供协作实例使用。也可以声明 \co@global@ 或 \co@local@ 事件，具体可见范围与约束由目标平台决定。

事件也可以声明为数组：

\begin{lstlisting}[language=croqtile,caption={声明事件数组}]
shared event e[4], e1;
\end{lstlisting}

逗号分隔的声明可以同时定义数组 \co@e@ 和单个事件 \co@e1@。

\co@trigger@ 将事件置为 \emph{SET}：

\begin{lstlisting}[language=croqtile,caption={触发事件数组}]
shared event e[4], e1;
trigger e, e1;
\end{lstlisting}

\co@trigger e, e1@ 触发数组 \co@e@ 的各个事件和单个事件 \co@e1@。

\co@wait@ 等待事件触发。事件为 \emph{UNSET} 时，当前实例阻塞；变为 \emph{SET} 后，等待结束。

下面用事件让一个异步区域等待另一个区域：

\begin{lstlisting}[language=croqtile,caption={等待事件}]
parallel p by 2 {
  shared event e;                          // init 'e' to 'UNSET'
  inthreads.async (p == 1) { wait e; }     // wait-thread
                                           // once unblocked, change 'e' to 'UNSET'
  inthreads.async (p == 0) { trigger e; }  // trigger-thread
                                           // convert 'e' to 'SET'
  sync.barrier : thread;                   // sync point
}
\end{lstlisting}

索引 \co@1@ 的区域先等待 \co@e@，索引 \co@0@ 的区域触发它后，等待才能结束。两个区域随后到达 \co@sync.barrier@ 同步点。事件明确了触发操作必须先于等待结束的顺序。

等待完成后，事件自动恢复为 \emph{UNSET}，可以用于下一次触发与等待。


\subsection*{多级流水线的事件数组}

设想生产者装入分块，消费者使用分块计算。读取槽位前，消费者需要知道数据已经就绪；覆盖槽位前，生产者需要知道上一次读取已经结束。将这两个信号命名为 \co@full@ 和 \co@empty@。若流水线包含多个可复用槽位，每个槽位就各自配备一对事件：

\begin{lstlisting}[language=croqtile,caption={多级流水线的事件数组}]
shared event full[STAGES], empty[STAGES];
\end{lstlisting}

每个缓冲槽对应一组 \co@full@ 和 \co@empty@ 事件。\co@stage = iv_k % STAGES@ 将循环索引映射到固定数量的槽位：生产者等待空槽，写完后触发 \co@full@；消费者等待数据就绪，读完后再触发 \co@empty@。这样，同一个缓冲区可以在后续迭代中复用。


\subsection*{初始化空槽事件}

循环开始前，需要先触发空槽对应的 \co@empty@ 事件：

\begin{lstlisting}[language=croqtile,caption={引导协议}]
foreach {s} in [STAGES] {
  trigger empty[s];
}
\end{lstlisting}

事件初始为未触发状态。如果没有这一步，生产者第一次等待空槽时就无法继续。初始化表示这些槽位已经可以使用。

完整的双缓冲和多阶段流水线示例见 3.8.5 节。先理解事件的触发与等待关系，再阅读这些组合用法。


\subsection*{事件的配对与复用}

使用事件时，应把每个 \co@wait@ 与相应的 \co@trigger@ 配对，并考虑等待完成后的自动复位。

每次等待都需要有可到达的触发操作；缓冲区循环复用时，还要安排好初始状态。



\section{分支与提前返回}
\label{sec:2-11}

\subsection*{运行时条件：\texttt{if}}

\co@if@ 根据运行时条件决定是否执行代码。它适合检查当前索引或运行值，再决定是否处理这份工作；没有进入分支的实例会继续执行分支后的语句：

\begin{lstlisting}[language=croqtile,caption={用 if 做谓词执行}]
if (tile_id < total_tiles) {
  // body executes only when the condition is true
}
\end{lstlisting}

条件为真的实例执行分支，条件为假的实例跳过。在按线程组执行的平台上，组内条件不一致可能让不同路径依次执行，增加执行成本。

运行时边界检查、条件累加等判断使用 \co@if@，条件可以依赖当前分块坐标或输入值。若语句体含有需要协作组共同参与的操作，应让组内实例对条件作出一致判断。

条件尽量在同一执行组内保持一致。本例所面向的线程束包含 32 个线程。


\subsection*{从 \texttt{parallel} 提前返回}

\co@yield@ 让当前并行实例提前结束，生成对应的 \co@return;@。与跳过一个 \co@if@ 分支不同，它会省去该实例的所有后续工作，适合已经完成任务或无需继续处理的实例：

\begin{lstlisting}[language=croqtile,caption={从 parallel 提前返回}]
parallel p by BLK_CNT: block {
  parallel q by THR_CNT : thread {
    if (cond) yield;
    ...
  }
}
\end{lstlisting}

\begin{keypoint}
\co@yield@ 不会自动检查提前退出是否影响同步。使用前应确认其他实例仍能完成相应的等待与同步操作。
\end{keypoint}

检查提前退出时，要对照区域中尚未执行的同步操作。实例若在触发事件前离开，消费者可能一直等待；若跳过必须共同参与的操作，也可能阻塞其他实例。负责同一分块的实例做出一致的尾部判断时，最容易确认这些同步责任是否仍能完成。



\section{四步 MMA 语法}
\label{sec:2-12}

矩阵乘法中，一个输出分块的结果通常来自多个 K 维切片。程序先准备累加器，再逐次装入对应的 A、B 分块，把每次乘加结果留在同一个累加器中，最后写回输出。MMA 把这个过程分为四步：初始化累加器、装入操作数、计算并累加、写回结果。

\begin{itemize}[leftmargin=1.6em]
\item \co@mma.fill 0.0@：初始化累加器。
\item \co@mma.load@：装入当前分块的操作数。
\item \co@mma.row.row mc, ma, mb@：执行 C += A × B，将结果累加到 \co@mc@。
\item \co@mma.store mc, dst@：写回结果。
\end{itemize}

装入与计算步骤沿 K 维重复，每次把一个切片的结果累加到 \co@mc@；完成后再写回输出。底层的线程数据分配由编译器结合布局约定处理。

\co@row.row@ 表示两个操作数采用行优先布局。其他组合包括 \co@row.col@、\co@col.row@ 和 \co@col.col@。布局必须与输入数据一致，否则计算结果可能错误。

若 K 被分为两个切片，先初始化一次 \co@mc@，装入第一对操作数并累加，再装入第二对并累加到同一个 \co@mc@，最后才写回。装入句柄描述操作数片段，累加器则跨迭代保留部分结果，因此填充与写回应放在反复执行的装入和乘加过程之外。



\section{C++ 嵌入、调用、宏与预处理器}
\label{sec:2-13}

\subsection*{\texttt{\_\_cpp\_\_}：原样注入 C++}

\co@__cpp__@ 将字符串作为 C++ 源码插入到它在 Croqtile 函数中的位置。某个简短目标操作需要放在特定并行或异步区域内时，可以使用这种形式。插入内容必须在该作用域中合法，引用的名字也必须对目标编译器可见。Croqtile 原样传递这些内容，不解析或改写字符串内部的代码。

\begin{itemize}[leftmargin=1.6em]
\item \co@__cpp__("...")@：适合简短代码。
\item \co@__cpp__(R"(...)")@：原始字符串写法，适合包含引号或多行内容的代码。
\end{itemize}


预处理在 Croqtile 解析之前进行。普通 C++ 和 \co@__co__@ 区域可以共用宏定义，便于统一维度与配置。

\begin{table}[htbp]\centering\small
\begin{tabularx}{\linewidth}{@{}lX@{}}
\toprule
指令 & 作用 \\
\midrule
\texttt{\#define NAME value} & 对象式宏：文本替换 \\
\texttt{\#if} / \texttt{\#elif} / \texttt{\#else} / \texttt{\#endif} & 条件包含 \\
\texttt{\#ifdef} / \texttt{\#ifndef} & 判断宏是否已定义的简写 \\
\texttt{\#error message} & 带消息强制编译时失败 \\
\bottomrule
\end{tabularx}
\end{table}

预处理器不支持函数式宏。需要带参数的计算时，可以在普通 C++ 中使用辅助函数。


\subsection*{集中定义分块参数}

分块尺寸与阶段数可以集中定义为宏：

\begin{lstlisting}[language=croqtile,caption={用宏表示分块参数}]
#define MATMUL_WARP_M 64
#define MATMUL_WARP_N 128
#define MATMUL_TILE_K 64
#define MATMUL_WARP_K 16
#define MATMUL_STAGES 4
\end{lstlisting}

这些名字可以在 \co@parallel@、缓冲区声明、循环范围和主机侧检查中共同使用。


\subsection*{用 \texttt{\#error} 做编译时断言}

\begin{lstlisting}[language=croqtile,caption={用 \#error 做编译时断言}]
#if MATMUL_SWIZ != (2 * MATMUL_TILE_K)
#error "MATMUL_SWIZ must equal 2 * MATMUL_TILE_K for f16 kernel"
#endif
#if MATMUL_WARP_M != 64
#endif
\end{lstlisting}

预处理器在解析 Croqtile 函数前求值这一条件。配置常量不兼容时，构建便会停止并显示给定消息。这适合检查配置宏之间的关系；若属性取决于运行时输入数据，则应使用运行时断言。


\subsection*{条件编译的代码路径}

\begin{lstlisting}[language=croqtile,caption={条件代码路径}]
#define PATH0
__co__ auto foo() {
  #ifdef PATH0
    // path 0 code
  #else
    // path 1 code
  #endif
}
\end{lstlisting}

预处理器根据条件保留相应的代码分支。宏也可以从命令行定义，例如 \co@croq kernel.co -DMATMUL_TILE_K=128@。


\subsection*{调用 C++ 辅助函数}

Croqtile 函数不是设备函数：在另一个函数里调用它不会产生设备调用，编译器会警告找不到同名的设备函数；调用 C++ 辅助函数才是常规用法。

\begin{lstlisting}[language=croqtile,caption={用 call 调用 C++ 辅助函数}]
void bar() {...}    // C++ kernel function
__co__ void foo() {
  parallel p by 6 {
    call bar();     // Call the C++ function
  }
}
\end{lstlisting}

本例用 \co@call bar()@ 调用 C++ 辅助函数 \co@bar@。设备调用的作用域和参数约束见本节后面的说明。


同一个配置常量经常同时用于主机分配和核函数声明。将它定义为对象式宏，可以让这些位置使用同一份配置。下面的例子说明，在 Croqtile 检查形状和类型之前，文本替换与条件编译怎样影响源码。


\subsection*{对象式宏}

参数化宏同样会展开，参数在形状求值前完成替换：

\begin{lstlisting}[language=croqtile,caption={对象式宏}]
#define M 256
#define N 32
#define K 64
__co__ auto matmul(f32 [M, K] lhs, f32 [N, K] rhs) { /*...*/ }
void foo() {
  croq::f32 a[M][K];
  croq::f32 b[N][K];
  // ...
  auto res = matmul(croq::make_spanview<2>(a, {M, K}),
                    croq::make_spanview<2>(b, {N, K}));
}
\end{lstlisting}

本例的 \co@M@、\co@N@、\co@K@ 在解析前被替换为 \co@256@、\co@32@、\co@64@，因此对应的形状是静态确定的。主机代码也使用相同的宏值。

带参数的函数式宏不受支持：

\begin{lstlisting}[language=croqtile,caption={带参数的宏}]
#define MAX(a, b) (((a) > (b)) ? (a) : (b))
\end{lstlisting}

这个 \co@MAX@ 定义不能按函数式宏使用。


\subsection*{注释}

注释可以使用 \co@/*...*/@ 或 \co@//...@。


\subsection*{条件编译}

条件编译支持 \co@#if@、\co@#ifdef@、\co@#ifndef@、\co@#else@ 和 \co@#endif@ 等指令：

\begin{lstlisting}[language=croqtile,caption={条件编译}]
#define PATH0
// some host code
__co__ auto foo() {
#ifdef PATH0
// some code related to PATH0
#else
// other code
#endif
}
// host control
#ifdef PATH0
// ...
#else
// ...
#endif
\end{lstlisting}

同一个条件可以同时控制主机代码与 Croqtile 代码。

这里使用 Croqtile 预处理器支持的指令即可；其支持范围不等同于完整的 C 预处理器。


设备辅助函数用于完成具体计算，Tileflow 程序负责在合适的执行位置调用它。


\subsection*{调用设备函数}

设备函数使用 \co@call@ 调用。调用需要放在可以执行设备代码的并行区域中：Tileflow 程序先准备数据并选择执行位置，再把相应参数交给设备函数。它主要用于按需复用已有硬件函数；硬件支持完善时，通常直接使用 Croqtile 提供的能力。


\subsection*{调用语句的基本语法}

\co@call@ 的基本形式如下：

\begin{lstlisting}[language=croqtile,caption={调用语句的一般语法}]
call func-name <optional-template-args> (arguments);
\end{lstlisting}

函数名后可以附带尖括号中的模板实参，普通参数写在圆括号中，并用逗号分隔。

目标编译器检查被调函数是否存在，以及其签名是否接受生成的实参。Croqtile 不在转译时执行这些 C++ 检查。因此，调用错误可能指向生成源码，需要对照辅助函数声明与原始调用位置查看。

\co@foo@ 中的调用位于并行区域内；\co@mou@ 中的调用处于主机执行区域，不能直接调用设备函数。

\begin{lstlisting}[language=croqtile,caption={非法的调用点}]
__co__ void foo() {
  parallel p by 1 {
    call bar();  // ok
  }
}
__co__ void mou() {
  call bar();  // error: not able to call device function from host
}
\end{lstlisting}


\subsection*{设备调用的参数类型}

设备调用可传递以下参数：

\begin{itemize}[leftmargin=1.6em]
\item 带形状数据。
\item 标量整数。
\item 目标平台支持的标量浮点值。
\end{itemize}

控制整数可以使用 \co@int@、\co@s32@ 或 \co@u32@。可传递的浮点类型取决于目标平台，例如 \co@float@、\co@double@、\co@half@、\co@bfp16@ 或 \co@half8@。

Tileflow 侧不进行这些浮点值的算术运算，而是把它们传给设备函数用于计算。

下面将几类参数传给设备辅助函数：

\begin{lstlisting}[language=croqtile,caption={向设备函数传递参数}]
__device__ void bar(float *p, int m, int n) {}
__device__ void foo(float *, int, unsigned, double, float) {}
__co__ void foobar(f32 [M, 24] input, int N, float padding) {
  parallel p by 1 {
    shared f32 [14, 7] buffer;
    call bar(input, input.span(0), buffer.span(0));
    call foo(buffer, N, 3, 3.14, padding);
  }
}
\end{lstlisting}

带形状数据传到设备函数时，参数是指向元素的指针，不再携带完整形状。例如 \co@bar@ 的 \co@float *p@ 对应 \co@input.data@；需要的维度另以整数参数传递。设备函数的声明应与这些参数类型匹配。

目标平台没有对应的原生类型时，可以使用 \co@croq::f16@、\co@croq::bf16@ 等类型。


\subsection*{设备调用的检查阶段}

设备调用的接口由后续目标编译阶段检查。出现参数或类型错误时，需要结合生成的调用和设备函数声明定位问题。



\section{打印、断言与编译时检视}
\label{sec:2-14}

\subsection*{\texttt{print!} 与 \texttt{println!}：编译时检视}

\co@print!@ 和 \co@println!@ 在编译时输出信息，可用于查看形状、范围和类型。例如，修改分块因子后，可以先打印推导出的形状，确认尺寸关系成立，再运行程序检查数值结果：

\begin{lstlisting}[language=croqtile,caption={编译时打印形状}]
__co__ void check_shapes(f32 [3, 2] b) {
  print!("shape of b: ");
  println!("b.span = ", b.span);
}
\end{lstlisting}

编译时得到如下输出：

\begin{lstlisting}[language=croqtile,caption={打印出的形状}]
shape of b: b.span = [3, 2]
\end{lstlisting}

相邻字符串字面量会拼接，例如 \co@"wor" "ld!"@ 等价于 \co@"world!"@。


\subsection*{\texttt{print} 与 \texttt{println}：运行时设备输出}

\begin{lstlisting}[language=croqtile,caption={一个运行时检视核函数}]
__co__ void inspect(s32 [4, 8] data) {
  foreach i in [data.span] {
    println("element ", i, " = ", data.at(i));
  }
}
\end{lstlisting}

运行时打印接受字符串和表达式，参数用逗号分隔。它适合观察实际执行中的索引或数据，与编译时检视互相补充。多个并行实例的输出顺序不固定，因此读取日志时需要结合实例索引辨认来源。

若只需检查一条执行路径，可以用索引条件限制打印：

\begin{lstlisting}[language=croqtile,caption={并行块内的打印}]
parallel {px, py} by [8, 16] : block
  parallel {qx, qy} by [16, 16] : thread {
    // ... compute ...
    if (px == 0 && py == 0 && qx == 3 && qy == 5) {
      println("partial sum = ", accum);
    }
  }
\end{lstlisting}

调试时可以用条件限制打印的实例，避免大量重复输出。


\subsection*{\texttt{assert}：运行时不变式检查}

\co@assert@ 在运行时检查条件，不满足时输出消息并中止。适合把程序必须满足的假设写在使用位置附近，让失败信息直接指向相应索引、数据或执行阶段：

\begin{lstlisting}[language=croqtile,caption={assert：运行时不变式检查}]
assert(stage < MATMUL_STAGES, "stage index out of bounds");
\end{lstlisting}

可以用断言检查索引范围、指针有效性或其他必须成立的条件，以便尽早发现错误。



\section{语法速查}
\label{sec:2-15}

\subsection*{数据搬运与分块选择}

下面汇总数据搬运和分块选择的常用写法。

\begin{table}[htbp]\centering\small
\begin{tabularx}{\linewidth}{@{}lX@{}}
\toprule
语法 & 含义 \\
\midrule
\texttt{dma.copy.async ...} & 启动非阻塞 DMA 传输。 \\
\texttt{wait f} & 等待 \texttt{f} 对应的传输完成。 \\
\texttt{future.data} & 访问已搬运分块。 \\
\texttt{future.span} & 访问已搬运分块的形状。 \\
\texttt{dma.transp\textless{}...\textgreater{}} & 搬运过程中置换维度。 \\
\texttt{dma.pad\textless{}...\textgreater{}} & 搬运过程中填充，并显式给出填充值。 \\
\texttt{.zfill} & 把不完整分块中的越界位置补零。 \\
\texttt{.swiz\textless{}N\textgreater{}} & 选择共享内存重排布局。 \\
\texttt{chunkat(...)} & 从规则分块中选出一个等分块。 \\
\texttt{subspan(...).at(...)} & 按显式分块范围与分块坐标选取分块。 \\
\texttt{view(...).from(...)} & 按显式范围与显式元素偏移选取分块。 \\
\texttt{tile \# i} & 把分块坐标与块内坐标合成为全局坐标。 \\
\texttt{\#tile} & 求 \texttt{tile} 轴的范围。 \\
\bottomrule
\end{tabularx}
\end{table}

阅读搬运语句时，先确定所选区域和目标位置，再检查目标何时就绪。流水线中仍需回答这两个问题；额外缓冲区只是让不同分块能够同时处在这个过程的不同阶段。


\subsection*{并行与迭代}

\begin{table}[htbp]\centering\small
\begin{tabularx}{\linewidth}{@{}lX@{}}
\toprule
\textbf{语法} & \textbf{含义} \\
\midrule
\texttt{parallel p by N} & N 个独立实例，索引为 0 至 N-1 \\
\texttt{parallel px, py by [M, N]} & 同一并行层次中的两个坐标轴 \\
\texttt{parallel p by A, q by B} & 逗号分隔的多个并行层次 \newline{}（\texttt{p} 与 \texttt{q} 对应不同并行层次） \\
\texttt{parallel q by N : thread} & 映射到线程块内的线程 \\
\texttt{foreach index = m, n, k in [a, b, c]} & 在 \texttt{foreach} 中对具名元组解构 \\
\texttt{lhs.span(0)} & 取出张量形状的一个维度 \\
\texttt{p \# m} & 把外部分块索引 \texttt{p} 与内部偏移 \texttt{m} 合成为全局坐标 \\
\bottomrule
\end{tabularx}
\end{table}

2.12 节概述 MMA 的计算步骤，3.7 节给出完整示例。


\subsection*{分支与协调}

\begin{table}[htbp]\centering\small
\begin{tabularx}{\linewidth}{@{}lX@{}}
\toprule
\textbf{语法} & \textbf{含义} \\
\midrule
\texttt{if (expr) ...} & 按运行时条件选择执行 \\
\texttt{inthreads.async (condition)} & 结构化并发区域：编译时划分线程 \\
\texttt{shared event name[N]} & 声明 N 个同步令牌 \\
\texttt{trigger name} & 发出条件已满足的信号 \\
\texttt{wait name} & 阻塞直到相应的 \texttt{trigger} 触发 \\
\texttt{tile\_id = tile\_iter \# block\_id} & 组合分块索引 \\
\texttt{int total\_tiles = expr} & 初始化一个控制整数 \\
\texttt{yield} & 从 \texttt{parallel} 块提前返回 \\
\bottomrule
\end{tabularx}
\end{table}


\subsection*{期值与循环简写}

\begin{table}[htbp]\centering\small
\begin{tabularx}{\linewidth}{@{}lX@{}}
\toprule
| \textbf{语法} & \textbf{含义} | \\
\midrule
| --- & --- | \\
| \co@dma.copy.async src => dst@ & 非阻塞拷贝，立即返回 | \\
| \co@dma.any@ & 占位期值，尚无传输在途 | \\
| \co@swap(f0, f1)@ & 交换两个期值句柄，不拷贝数据 | \\
| \co@rotate(f0, f1, f2)@ & 轮转三个期值句柄 | \\
| \co@with tile_k in N ...@ & 绑定到范围为 N 的分块轴，作用域受限 | \\
| \co@foreach tile_k(1:)@ & 从索引 1 开始迭代 | \\
| \co@__co__ auto fn(...)@ & 由 \co@return@ 语句推导返回类型 | \\
\bottomrule
\end{tabularx}
\end{table}



\guidechapternext{You can now read most \file{.co} files. Chapter~\ref{ch:programming-guide} explains why the constructs are shaped this way and covers the advanced forms; Appendix~\ref{app:quickref} is a compact lookup table.}

